% Part 4 bench circuits, redrawn from the handout (steps 8 and 14), with
% the DMM drawn as its internal resistance R_VM in a dashed box (step 15).
% Component names follow the handout: S1, R8, R10, C13, V5 for charging;
% S2, R11, R12, C14, V6 for discharging. The handout's charging drawing says
% 16.5 V; its note corrects that to 15 V.
% spdt with rotate=90: pole (in) at the bottom, blade to the left throw
% (out 1). Adding yscale=-1 moves the blade to the right throw.
\begin{circuitikz}[x=1cm, y=1cm, font=\small,
    every label/.append style={font=\small}]
  \ctikzset{bipoles/length=0.9cm, resistors/scale=0.8, capacitors/scale=0.8, sources/scale=1.2}
  \tikzset{head/.style={font=\sffamily\bfseries\fontsize{8.5}{10}\selectfont,
      anchor=base west, inner sep=0pt},
    sub/.style={font=\sffamily\fontsize{7.5}{9}\selectfont, text=chartText,
      anchor=base west, inner sep=0pt},
    meter/.style={dashed, gray, line width=0.4pt}}
  % ---------- Charging (step 8) ----------
  \begin{scope}[shift={(0,0)}]
    \node[head] at (0,4.0) {Charging circuit \textmd{\textperiodcentered\ step 8}};
    \node[spdt, rotate=90] (S1) at (2.9,2.35) {};
    \node[font=\small, anchor=west] at ([xshift=2pt]S1.out 2) {S1};
    \draw (0,0) to[V, invert, l2_={V5 and \qty{15}{\volt}}, l2 valign=c] (0,3.1)
          to[R, l={R8}, a={\qty{1}{\mega\ohm}}] (S1.out 1 |- 0,3.1) -- (S1.out 1);
    \draw (S1.out 2) |- (3.15,3.1) -- (6.1,3.1)
          to[R, l2^={R10 and \qty{47}{\ohm}}, l2 valign=c] (6.1,0);
    \draw (S1.in) -- (2.9,1.6) node[circ] {}
          to[cC, l2_={C13 and \qty{47}{\micro\farad}}, l2 valign=c] (2.9,0);
    \draw (2.9,1.6) -- (4.3,1.6)
          to[R, l2^={\(R_{\mathrm{VM}}\) and \qty{9.6}{\mega\ohm}}, l2 valign=c] (4.3,0);
    \draw[meter] (3.75,-0.25) rectangle (5.72,1.95);
    \node[sub, anchor=south east] at (5.72,1.98) {XMM1 (DMM)};
    \draw (0,0) -- (6.1,0);
    \draw (2.9,0) node[circ] {} (4.3,0) node[circ] {} (0,0) node[circ] {} node[ground] {};
  \end{scope}
  % ---------- Discharging (step 14) ----------
  \begin{scope}[shift={(8.35,0)}]
    \node[head] at (0,4.0) {Discharging circuit \textmd{\textperiodcentered\ step 14}};
    \node[spdt, rotate=90, yscale=-1] (S2) at (2.9,2.35) {};
    \node[font=\small, anchor=west] at ([xshift=2pt]S2.out 1) {S2};
    \draw (0,0) to[V, invert, l2_={V6 and \qty{15}{\volt}}, l2 valign=c] (0,3.1)
          to[R, l={R11}, a={\qty{47}{\ohm}}] (S2.out 2 |- 0,3.1) -- (S2.out 2);
    \draw (S2.out 1) |- (3.15,3.1) -- (6.1,3.1)
          to[R, l2^={R12 and \qty{1}{\mega\ohm}}, l2 valign=c] (6.1,0);
    \draw (S2.in) -- (2.9,1.6) node[circ] {}
          to[cC, l2_={C14 and \qty{47}{\micro\farad}}, l2 valign=c] (2.9,0);
    \draw (2.9,1.6) -- (4.3,1.6)
          to[R, l2^={\(R_{\mathrm{VM}}\) and \qty{9.6}{\mega\ohm}}, l2 valign=c] (4.3,0);
    \draw[meter] (3.75,-0.25) rectangle (5.72,1.95);
    \node[sub, anchor=south east] at (5.72,1.98) {XMM2 (DMM)};
    \draw (0,0) -- (6.1,0);
    \draw (2.9,0) node[circ] {} (4.3,0) node[circ] {} (0,0) node[circ] {} node[ground] {};
  \end{scope}
\end{circuitikz}
